\chapter{Compositional Safety for Battery Fleets}
\label{ch:compositional-safety-fleets}

This chapter keeps the thesis inside the battery domain while pushing beyond
the single-asset case.  The current locked extension is a two-battery shared
transformer study, packaged through
\texttt{reports/fleet/fleet\_composition\_two\_battery.csv} and
\texttt{reports/fleet/two\_battery\_composition\_metrics.csv}.

\section{Why Composition Matters}

Two individually safe batteries can still create a fleet-level problem if
they respond to local conditions without coordinating at the point of
interconnection.  In battery terms, the risk is not only a local SOC breach;
it is also aggregate export or curtailment pressure on the shared transformer.

\section{Battery-Fleet View}

The fleet extension keeps the ORIUS battery logic local to each unit:

\begin{enumerate}
  \item Each battery produces its own candidate safe action under degraded
    telemetry.
  \item The shared transformer then enforces a fleet-level envelope on the
    combined dispatch.
  \item Any excess is absorbed as coordinated curtailment rather than as an
    unsafe aggregate command.
\end{enumerate}

This is the battery-operational meaning of composition: preserve local safety
while making sure the shared export path never becomes the new hidden gap.

\section{Locked Two-Battery Stress Output}

The present wrapper-generated fleet metrics report:

\begin{itemize}
  \item 56 joint over-limit steps,
  \item a maximum curtailment of 10.0\,MW, and
  \item 4513.97\,MWh of useful work preserved.
\end{itemize}

These values come directly from
\texttt{reports/fleet/two\_battery\_composition\_metrics.csv}.

\begin{table}[ht]
\centering
\caption{Locked fleet-composition metrics from the current two-battery study.}
\label{tab:fleet-results}
\begin{tabular}{lrr}
\toprule
\textbf{Metric} & \textbf{Value} & \textbf{Battery consequence} \\
\midrule
Joint over-limit steps & 56 & Shared export stress appears repeatedly \\
Max curtailment (MW) & 10.0 & Coordination trims peak aggregate dispatch \\
Useful work preserved (MWh) & 4513.97 & Safety does not force total shutdown \\
\bottomrule
\end{tabular}
\end{table}

\section{Multi-Agent Protocol Comparison}

Table~\ref{tab:multi_agent} reports the non-composition counterexample
quantitatively: two batteries each proposing 60~MW of discharge into a
shared 80~MW feeder.  The independent protocol produces 24 joint
violations; ORIUS-coordinated (centralized) produces zero, while
preserving the same useful work output.

\begin{table}[htbp]
\centering
\small
\caption{Multi-agent coordination under a shared 80\,MW transformer limit. ORIUS eliminates joint violations without reducing useful work.}
\label{tab:multi_agent}
\begin{tabular}{@{}lrrrr@{}}
\toprule
Protocol & Joint viol. & Local viol. & Work (MWh) & Margin \\
\midrule
Independent & 24 & 0 & 1920 & 0 \\
ORIUS coordinated & \textbf{0} & 0 & 1920 & 1 \\
Distributed & \textbf{0} & 0 & 1920 & 1 \\
\bottomrule
\end{tabular}
\end{table}

\section{Interpretation}

The important battery result is not that the fleet chapter is ``finished'' as
a full theorem-backed composition theory.  It is that the repo now contains a
locked, chapter-facing composition run showing exactly where shared
infrastructure stress emerges and how much useful work remains after
coordination.  That closes a practical gap between single-battery safety and
fleet operation without pretending that all composition questions are solved.

\section{What Remains Open}

The present fleet package is intentionally narrow:

\begin{itemize}
  \item It is a two-battery study, not a heterogeneous multi-asset proof.
  \item It reports aggregate stress and preserved work, not a full fairness or
    degradation-allocation frontier.
  \item It motivates the quality-score gossip and coordinated fallback ideas,
    but does not yet prove their optimality.
\end{itemize}

\section{Summary}

The battery thesis now includes a real fleet-composition artifact rather than
only a conceptual extension.  The locked two-battery study shows repeated
shared-transformer stress, bounded curtailment, and non-trivial useful work
preservation.  For the battery book, that is the correct claim: composition
has moved from future-work rhetoric into an auditable battery extension, while
still remaining narrower than a full universal coordination theory.
